% !TEX root = ../main.tex
\section{Introduction}\label{sec:introduction}

Barkeshli, Douglas, and Freedman~\cite{BDF2026} propose that mathematical reasoning carries a directed hypergraph structure, and illustrate this framework using CIC expressions, but do not instantiate it on a concrete proof system.  We construct a CIC kernel that performs type-checking directly on a content-addressed hypergraph (where each node's identity is determined by the hash of its content).  The kernel's sole data structure is a content-addressed hypergraph.  Each step of type-checking takes the hypergraph as input and returns a larger hypergraph as output.  We implement a traditional expression-tree kernel in Lean~4 as a baseline and verify that both kernels behave identically when checking \texttt{Nat.add\_comm}.

This construction yields two immediate results.  The output of type-checking is a complete hypergraph containing all declarations, sub-expressions, and intermediate expressions produced during verification, amenable to direct analysis using the depth functions and inductive closure of the BDF framework.  Content-addressing ensures that each node's identity is determined by its content hash, independent of any naming system or expression representation.

Our work draws on three lines of research.  On the structural side, Barkeshli, Douglas, and Freedman~\cite{BDF2026} characterize the hypergraph structure of mathematical reasoning from a theoretical perspective.  Aksenov, Bodnia, Freedman, and Mulligan~\cite{Aksenov2026} propose that the compressibility of mathematical knowledge reflects its depth.  Li, Peng, Shafto, and Severini~\cite{LPSS26b} analyze the multinetwork structure of Mathlib from an empirical perspective.  On the content-addressing side, Lean~4 optimizes performance at the C++ level through hash consing and sub-expression sharing, but these remain invisible to the logical specification.  Yatima~\cite{Arg23} content-addresses Lean~4 expressions to support zero-knowledge verification.  zkPi~\cite{zkPi} places a CIC kernel inside a zkSNARK circuit.  Unison uses content-addressing for a general-purpose programming language.  On the kernel implementation side, lean4lean~\cite{lean4lean} re-implements Lean's kernel in Lean~4.  In all these systems, type-checking operates on expression trees.  Our construction unifies hypergraph structure and content-addressing at the kernel execution level.
